\section*{第 \ref{ch:b3:exploration-seismology} 章——勘探地震学}

\begin{solution}{exo:b3:exploration-seismology:1}
$Z = \rho v$：白垩 $2000\times2500 = \qty{5.00e6}{kg.m^{-2}.s^{-1}}$，页岩
$2300\times2300 = \qty{5.29e6}{kg.m^{-2}.s^{-1}}$，石灰岩 $2600\times4500 =
\qty{11.70e6}{kg.m^{-2}.s^{-1}}$。白垩--页岩界面：$R = (5.29 - 5.00)/(5.29 + 5.00)
= 0.29/10.29 = 0.0282$；页岩--石灰岩界面：$R = (11.70 - 5.29)/(11.70 + 5.29) =
6.41/16.99 = 0.377$。反射能量为 $R^2 = 0.142$：石灰岩顶面
送回入射能量的 \qty{14.2}{\percent}，白垩–页岩
界面不足 \qty{0.1}{\percent}。
\end{solution}

\begin{solution}{exo:b3:exploration-seismology:2}
$t_0 = 2\times800/2000 = \qty{0.800}{s}$；$t = \sqrt{0.800^2 + (1000/2000)^2} =
\sqrt{0.89} = \qty{0.9434}{s}$；时差为 \qty{0.1434}{s}。
近似式给出 $1000^2/(2\times2000^2\times0.800) = \qty{0.1563}{s}$，偏大
\qty{9}{\percent}：炮检距大于深度时，$x \ll vt_0$ 不再
成立。
\end{solution}

\begin{solution}{exo:b3:exploration-seismology:3}
中点的移动量为检波器间距的一半：\omterm{def:b3:exploration-seismology:cmp}{共中心点}间距为
$12.5/2 = \qty{6.25}{m}$。\omterm{def:b3:exploration-seismology:cmp}{覆盖次数}为 $240\times12.5/(2\times25) = 60$。若要
\omterm{def:b3:exploration-seismology:cmp}{覆盖次数}达到 120，炮点间距需为 $240\times12.5/(2\times120) = \qty{12.5}{m}$。
\end{solution}

\begin{solution}{exo:b3:exploration-seismology:4}
$\lambda = v/f$：砂岩中 \qty{100}{m}，$\lambda/4 = \qty{25}{m}$；
石灰岩中 \qty{150}{m}，$\lambda/4 = \qty{37.5}{m}$。取
$\lambda = 2800/30 = \qty{93.3}{m}$ 与 $z = \qty{1000}{m}$，
$r_F \approx \sqrt{\lambda z/2} = \sqrt{46\,667} = \qty{216}{m}$（\qty{217}{m}
计入 $\lambda^2/16$ 项后）。
\end{solution}

\begin{solution}{exo:b3:exploration-seismology:5}
\omterm{thm:b3:exploration-seismology:dix}{迪克斯公式}：$v^2 = (2800^2\times0.700 - 2400^2\times0.500)/0.200 = (5\,488\,000 -
2\,880\,000)/0.200 = \qty{1.304e7}{m^2.s^{-2}}$，故 $v = \qty{3611}{m.s^{-1}}$。
单程穿越该层所需时间为 \qty{0.100}{s}：$h = 3611\times0.100 =
\qty{361}{m}$。
\end{solution}

\begin{solution}{exo:b3:exploration-seismology:6}
(a) 设 $d$ 为中点 M 到反射面的距离，于是
$t_0 = 2d/v$。自震源 S 到检波点 G 的反射射线
长度等于反射面中 S 的镜像 S$'$ 到 G 的直线段。取 M 为
原点，直线沿倾向；S$'$ 位于距 S 为 $2d_S$ 的位置，沿
法线方向，其中 $d_S = d - (x/2)\sin\varphi$。S$'$G 的平方为
$x^2 + 4d_S^2 + 4d_Sx\sin\varphi = x^2 + 4d^2 - x^2\sin^2\varphi$，所有
含 $dx$ 的项相消，于是 $v^2t^2 = 4d^2 + x^2\cos^2\varphi$，且
$t^2 = t_0^2 + x^2\cos^2\varphi/v^2$。(b) 道集服从速度为
$V = v/\cos\varphi = 2000/\cos\ang{12} = \qty{2045}{m.s^{-1}}$ 的双曲线：倾角使
\omterm{def:b3:exploration-seismology:velocities}{叠加速度}偏大，偏大幅度为 \qty{2.2}{\percent}。
\end{solution}

\begin{solution}{exo:b3:exploration-seismology:7}
由 $\sqrt{60} = 7.75$ 得。信噪比加倍需要
道数为原来的四倍：\omterm{def:b3:exploration-seismology:cmp}{覆盖次数}为 240。
\end{solution}

\begin{solution}{exo:b3:exploration-seismology:8}
绘制深度 $vt/2$ 每千米的变化量为 $\qty{1.5}{km.s^{-1}}\times\qty{0.40}{s}
= \qty{0.60}{km}$：$\tan\theta_a = 0.60$，$\theta_a = \ang{31.0}$。
由零炮检距定理，$\sin\theta = 0.60$：$\theta = \ang{36.9}$。
\end{solution}

\begin{solution}{exo:b3:exploration-seismology:9}
每个尖峰处放置一个以其为中心的 $w$ 副本：样点 0--2 上的 $0.3\,w$ 给出
$(-0.15, 0.3, -0.15)$，样点 2--4 上的 $-0.2\,w$ 给出 $(0.1, -0.2, 0.1)$。
二者之和为 $(0, -0.15, 0.3, -0.05, -0.2, 0.1, 0)$，从序列前一个样点
写到序列后一个样点。反射在样点 1 和 3 处达到峰值，
符号均正确；二者在样点 2 处重叠，此处 $-0.15 + 0.1 = -0.05$，
有轻微干涉，但仍是两个分开的峰。
\end{solution}

\begin{solution}{exo:b3:exploration-seismology:10}
$\arcsin(1800/2900) = \ang{38.37}$ 与 $\arcsin(1800/4400) = \ang{24.15}$，故
$i_c = (38.37 + 24.15)/2 = \ang{31.26}$，且 $\theta = (38.37 - 24.15)/2 =
\ang{7.11}$；$v_2 = 1800/\sin\ang{31.26} = \qty{3469}{m.s^{-1}}$。调和
平均 $2\times2900\times4400/(2900 + 4400) = \qty{3496}{m.s^{-1}}$，
高 \qty{0.8}{\percent}；乘以 $\cos\theta = 0.9923$ 后得到
\qty{3469}{m.s^{-1}}，即通常的小倾角近似
$v_2 \approx \cos\theta\,2v_dv_u/(v_d + v_u)$。
\end{solution}

\begin{solution}{exo:b3:exploration-seismology:11}
(a) 合成记录为 $R[w(t) - w(t - \Delta t)] \approx R\,\Delta t\,w'(t)$，由
一阶泰勒展开。(b) 其振幅正比于
$\Delta t$，因此比\omterm{def:b3:exploration-seismology:resolution}{调谐厚度}更薄的岩层仍会返回信号，
具有 $w'$ 的形状，但其顶和底不能拾取为两个事件：它
被探测到，但未被分辨。若 $R$ 与子波已知，例如来自
钻孔，则振幅可给出 $\Delta t$，进而给出厚度。(c) 振幅
正比于厚度：比值为 $5/10 = 1/2$。
\end{solution}

\begin{solution}{exo:b3:exploration-seismology:12}
(a) $t(x) = (2/v)\sqrt{z^2 + x^2}$：$t(0) = 2\times600/3000 = \qty{0.400}{s}$ 与
$t(300) = (2/3000)\sqrt{600^2 + 300^2} = \qty{0.4472}{s}$。(b) $t^2 = t_0^2 +
x^2/(v/2)^2$，即速度为 $v/2$ 的 NMO 形式。顶点时间相同的反射
在零炮检距剖面上完全不显示时差；与道集中同一 $t_0$ 的 NMO 双曲线相比，
绕射的时差
$x^2/(2(v/2)^2t_0)$ 大四倍，因为时间依赖于
散射体的距离两次，即下行与上行。
(c) 偏移沿此双曲线对剖面求和，把能量归于其顶点：
绕射收缩为位于 $x = 0$、$z = \qty{600}{m}$ 处的一点。
\end{solution}

\begin{solution}{pb:b3:exploration-seismology:1}
\textbf{1.} 由 W 端，$400/(486.1 - 425.2) = 400/0.0609 = \qty{6568}{m.s^{-1}}$；
由 E 端，$400/(486.1 - 372.4) = 400/0.1137 = \qty{3518}{m.s^{-1}}$。由 W 端的首波
速度似乎偏快：W 为下倾端，其炮点沿上倾方向传播；折射面
自 W 向 E 上升。

\textbf{2.} $\arcsin(2400/3518) = \ang{43.02}$，$\arcsin(2400/6568) =
\ang{21.43}$：$i_c = \ang{32.22}$，$\theta = \ang{10.79}$。

\textbf{3.} $v_2 = 2400/\sin\ang{32.22} = \qty{4501}{m.s^{-1}}$：处于
石灰岩的速度范围内，即侏罗纪石灰岩的顶部。

\textbf{4.} 截距为 \qty{1200}{m} 处的 $t - x/v_{\text{app}}$：
$t_{i,W} = 0.4861 - 1200/6568 = \qty{0.3034}{s}$，$t_{i,E} = 0.4861 -
1200/3518 = \qty{0.1450}{s}$。取 $\cos i_c = 0.8460$，$h = v_1t_i/(2\cos
i_c)$：$h_W = 2400\times0.3034/1.692 = \qty{430.4}{m}$，$h_E =
\qty{205.7}{m}$。

\textbf{5.} 除以 $\cos\theta = 0.9823$：W 下方为 \qty{438.1}{m}，
E 下方为 \qty{209.4}{m}。于是 $(438.1 - 209.4)/1200 = 0.1906 =
\tan\ang{10.79}$，即第 2 题的倾角。

\textbf{6.} 页岩不产生首波：它不可见，本次测量将其与白垩合并计算。
此处 $v_1 = \qty{2.40}{km.s^{-1}}$ 为两者的平均值；若整个覆盖层都采用白垩的直达波速度
\qty{2.50}{km.s^{-1}}，
则各深度都会偏大。

\textbf{7.} $V = 600/\sqrt{0.5437^2 - 0.4878^2} = 600/0.2401 =
\qty{2499}{m.s^{-1}}$，$V = 600/\sqrt{0.6240^2 - 0.5748^2} =
\qty{2471}{m.s^{-1}}$。第二个值将白垩与其下较慢的页岩一并平均，
故平均速度降低。

\textbf{8.} \qty{2811}{m.s^{-1}} 与 \qty{2829}{m.s^{-1}}。

\textbf{9.} 第 1 层：\qty{2499}{m.s^{-1}}。第 2 层：$v^2 = (2471^2\times0.5748
- 2499^2\times0.4878)/0.0870$，$v = \qty{2308}{m.s^{-1}}$。第 3 层：
$v^2 = (2811^2\times0.6548 - 2471^2\times0.5748)/0.0800$，
$v = \qty{4561}{m.s^{-1}}$。第 4 层：$v^2 = (2829^2\times0.7348 -
2811^2\times0.6548)/0.0800$，$v = \qty{2972}{m.s^{-1}}$。

\textbf{10.} 白垩（\qty{2499}{m.s^{-1}}）、页岩（\qty{2308}{m.s^{-1}}），
石灰岩（\qty{4561}{m.s^{-1}}）、砂岩（\qty{2972}{m.s^{-1}}），均处于
各自公布的范围之内。反射来自白垩与页岩的界面、石灰岩
顶面、石灰岩与砂岩的界面，以及
片岩基底顶部的不整合面。

\textbf{11.} 白垩与页岩的阻抗几乎相同：$R = 0.028$
（练习 1），而石灰岩顶面的为 0.38。

\textbf{12.} 速度为 $V_{\text{rms}}$ 的双曲线仅在偏移距的二阶近似下成立；在偏移距
\qty{2}{km}（为最深反射层深度的两倍）处，走时已偏离该双曲线，
若用整个排列拟合双曲线，
将得到错误的速度，迪克斯\omterm{def:b3:exploration-seismology:velocities}{层速度}则更甚。远偏移距
在校正时还会拉伸子波。

\textbf{13.} 各层的双程走时：\qty{0.4878}{s}、\qty{0.0870}{s}、
\qty{0.0800}{s}、\qty{0.0800}{s}。厚度 $v\,\Delta t/2$：\qty{609.5}{m}、
\qty{100.4}{m}、\qty{182.4}{m}、\qty{118.9}{m}。

\textbf{14.} 地面以下的深度：\qty{609.5}{m}、\qty{709.9}{m}、
\qty{892.3}{m}、\qty{1011.2}{m}。页岩、石灰岩与砂岩的结果分别与其厚度相差
不超过 \qty{0.4}{m}、\qty{2.4}{m} 与 \qty{1.1}{m}，误差在快速而薄的石灰岩中
增大，
此处微小的时间误差代价最大。

\textbf{15.} $\bar v = 2\times1011.2/0.7348 = \qty{2752}{m.s^{-1}}$。\omterm{def:b3:exploration-seismology:velocities}{均方根
速度}以 $v^2$ 对各层加权，且从不小于按走时加权的平均值：
$2829\times0.7348/2 = \qty{1039}{m}$，偏深 \qty{28}{m}。

\textbf{16.} $\lambda = 2308/30 = \qty{76.9}{m}$，$\lambda/4 = \qty{19.2}{m}$；
\omterm{def:b3:exploration-seismology:wavelet}{里克子波}的估计为 $2308/(5.2\times30) = \qty{14.8}{m}$。页岩厚
\qty{100}{m}，为其\omterm{def:b3:exploration-seismology:resolution}{调谐厚度}的五倍，其顶、底均可分辨。
按四分之一波长准则，它在 $f = v/(4b) =
2308/400 = \qty{5.77}{Hz}$ 处调谐。

\textbf{17.} $\lambda = 2829/30 = \qty{94.3}{m}$，$z = \qty{1011}{m}$：
$r_F \approx \sqrt{94.3\times1011/2} = \qty{218}{m}$（计入 $\lambda^2/16$ 项后为
\qty{220}{m}）。

\textbf{18.} $t(x)^2 = t_0^2 + (2x/v)^2$，故 $v = 2x/\sqrt{t^2 - t_0^2} =
400/\sqrt{0.5732^2 - 0.550^2} = 400/0.1614 = \qty{2478}{m.s^{-1}}$。

\textbf{19.} $z = vt_0/2 = 2478\times0.550/2 = \qty{681}{m}$。

\textbf{20.} 所绘深度 $vt/2$ 每单位长度的变化为 $1.24\times0.167 = 0.207$，
即 $\theta_a = \arctan0.207 = \ang{11.70}$。真倾角满足 $\sin\theta = 0.207$：
$\theta = \ang{11.95}$，即褶皱翼部的倾角。

\textbf{21.} $d = 2480\times0.655/2 = \qty{812}{m}$；偏移将反射点沿倾向向上移动
$d\sin\theta = 812\times0.207 = \qty{168}{m}$，朝向背斜顶部，
并使深度由 \qty{812}{m} 变为 $d\cos\theta =
\qty{795}{m}$。

\textbf{22.} $2500\times0.160/2 = \qty{200}{m}$。

\textbf{23.} 由第 13 题与第 14 题可知，在 CMP S 处，基底不整合
位于地面以下 \textbf{\qty[text-series-to-math=true]{1011}{m}} 处。
\end{solution}
